\documentclass[11pt,a4paper]{article}

\usepackage[T1]{fontenc}
\usepackage[utf8]{inputenc}
\usepackage{lmodern}
\usepackage{amssymb,amsmath,amsthm,mathtools}
\usepackage[margin=2.5cm]{geometry}
\usepackage{parskip}
\usepackage[hidelinks]{hyperref}
\usepackage{microtype}

% Theorem environments (arXiv / journal style)
\theoremstyle{plain}
\newtheorem{theorem}{Theorem}[section]
\newtheorem{proposition}[theorem]{Proposition}
\newtheorem{lemma}[theorem]{Lemma}
\newtheorem{corollary}[theorem]{Corollary}

\theoremstyle{definition}
\newtheorem{definition}[theorem]{Definition}
\newtheorem{example}[theorem]{Example}
\newtheorem{remark}[theorem]{Remark}

% Notation
\newcommand{\Ext}{\mathrm{Ext}}
\newcommand{\Shape}{\mathrm{Shape}}
\newcommand{\Pos}{\mathrm{Pos}}
\newcommand{\Type}{\mathbf{Type}}
\newcommand{\id}{\mathrm{id}}
\newcommand{\rootpos}{\mathsf{o}}
\newcommand{\sub}{\mathbin{\downarrow}}
\newcommand{\shift}{\oplus}
\newcommand{\cat}{\mathcal{C}}
\newcommand{\DCont}{\mathbf{DCont}}
\newcommand{\Cof}{\mathbf{Cof}}

\title{Morphisms of Directed Containers\\
  as Cofunctors}
\author{Extracted from a machine-checked development}
\date{}

\begin{document}
\maketitle

\begin{abstract}
A directed container presents a small category whose objects are shapes and
whose morphisms are positions. Morphisms of directed containers---shape maps
forward, position maps \emph{backward}, compatible with root, sub-shape, and
shift---are not functors of the presented categories: the variance on
hom-sets is opposite. They are precisely \emph{cofunctors} (retrofunctors).
We show that the two structures carry identical data, so there is a bijection
on hom-sets between directed-container morphisms and cofunctors; that
cofunctors form a strict category; and that the correspondence preserves
identities and composition, yielding an isomorphism of categories
$\DCont\cong\Cof$. The proofs of all stated results have been machine-checked;
we record only the mathematical content.
\end{abstract}

\tableofcontents
\bigskip

%======================================================================
\section{Introduction}
%======================================================================

A directed container equips a container $(S,P)$ with a root $\rootpos(s)$, a
sub-shape $s\sub p$, and a shift $p\shift q$, subject to laws D1--D5. Reading
these operations categorically, shapes are objects and positions out of a
shape are morphisms; the root is the identity and the shift is composition.
This is the object-level dictionary of Ahman--Uustalu: directed containers
present small categories.

At the morphism level a tension appears. A container morphism has a
\emph{covariant} shape map and a \emph{contravariant} position map. A functor
between the presented categories would act covariantly on hom-sets. The two
variances are opposite, so a directed-container morphism cannot be a functor.
With compatibility laws for root, sub-shape, and shift, it is exactly a
\textbf{cofunctor}: an object map together with a backward lift of morphisms.

\paragraph{What this note claims.}
In ordinary mathematical language: the data of a directed-container morphism
and the data of a cofunctor between the presented categories are the same
data, under the dictionary that renames root/sub/shift to identity/codomain/
composition. Consequently there is a hom-set bijection; cofunctors compose
strictly; and the bijection is an isomorphism of categories
$\DCont\cong\Cof$.

%======================================================================
\section{Directed containers and the presented category}
%======================================================================

We work in a single universe $\Type$. A \textbf{directed container} $D$ is a
container $(\Shape(D),\Pos(D))$ with root, sub-shape, and shift operations
satisfying the five directed laws D1--D5 (as in the companion note on directed
containers and the induced comonad). We write $S$ and $P$ for shapes and
positions when convenient.

Through the categorical dictionary we introduce the following vocabulary.

\begin{definition}[Dictionary]\label{def:dictionary}
Let $D$ be a directed container.
\begin{enumerate}
  \item The \textbf{identity morphism} at a shape $s$ is the root:
        $\id_s := \rootpos(s)\in P(s)$.
  \item The \textbf{codomain} of a position $p\in P(s)$ is the sub-shape:
        $\mathrm{cod}(s,p) := s\sub p\in S$.
  \item The \textbf{composition} of $p\in P(s)$ with
        $q\in P(s\sub p)$ is the shift:
        $q\circ p := p\shift q\in P(s)$
        (diagrammatic order: travel along $p$, then along $q$).
\end{enumerate}
\end{definition}

\begin{remark}
In this note we treat the presented category through this dictionary without
re-deriving the category laws; those laws are D2/D3 (units) and D5
(associativity), transported along D1 and D4, and are established in the
companion development on directed containers as categories.
\end{remark}

%======================================================================
\section{Morphisms of directed containers}
%======================================================================

\begin{definition}[Directed-container morphism]\label{def:dcont-morph}
Let $C$ and $D$ be directed containers. A \textbf{morphism of directed
containers} $\varphi:C\Rightarrow D$ consists of
\begin{enumerate}
  \item a map on shapes $f:\Shape(C)\to\Shape(D)$, and
  \item for each $s\in\Shape(C)$, a function
        $f^\sharp_s:\Pos(D)\bigl(f(s)\bigr)\to\Pos(C)(s)$
        (contravariant on positions),
\end{enumerate}
subject to the three compatibility laws M0--M2 below.
\end{definition}

\begin{definition}[Laws M0--M2]\label{def:m-laws}
Let $\varphi=(f,f^\sharp):C\Rightarrow D$ be as above.
\begin{enumerate}
  \item[\textbf{M0} (codomain).]
        For all $s\in\Shape(C)$ and $h\in\Pos(D)\bigl(f(s)\bigr)$,
        \[
          f\bigl(s\sub f^\sharp_s(h)\bigr)
          \;=\;
          f(s)\sub h.
        \]
  \item[\textbf{M1} (unit).]
        For all $s\in\Shape(C)$,
        \[
          f^\sharp_s\bigl(\rootpos(f(s))\bigr)
          \;=\;
          \rootpos(s).
        \]
  \item[\textbf{M2} (composition).]
        For all $s\in\Shape(C)$, $h\in\Pos(D)\bigl(f(s)\bigr)$, and
        $k\in\Pos(D)\bigl(f(s)\sub h\bigr)$,
        \[
          f^\sharp_s\bigl(h\shift k\bigr)
          \;=\;
          f^\sharp_s(h)
          \;\shift\;
          f^\sharp_{s\sub f^\sharp_s(h)}
            \bigl((\mathrm{M0}_{s,h})^{-1}\cdot k\bigr),
        \]
        where the inner position $k$ is transported along M0 so that the shift
        in $C$ is well-typed.
\end{enumerate}
\end{definition}

\begin{remark}
Law M2 transports along M0 exactly as the directed law D5 transports along
D4: the right factor of the shift lives in a fibre identified only up to a
shape equality.
\end{remark}

%======================================================================
\section{Cofunctors}
%======================================================================

\begin{definition}[Cofunctor]\label{def:cofunctor}
Let $C$ and $D$ be directed containers, and write $\cat(C)$ and $\cat(D)$ for
the categories they present. A \textbf{cofunctor} $\phi:\cat(C)\to\cat(D)$
consists of
\begin{enumerate}
  \item an object map $\mathrm{obj}:\Shape(C)\to\Shape(D)$, and
  \item for each $s\in\Shape(C)$, a \textbf{lift}
        $\mathrm{lift}_s:\Pos(D)\bigl(\mathrm{obj}(s)\bigr)\to\Pos(C)(s)$
        sending a morphism out of $\mathrm{obj}(s)$ to a morphism out of $s$,
\end{enumerate}
subject to the three cofunctor laws C0--C2 below.
\end{definition}

\begin{definition}[Laws C0--C2]\label{def:c-laws}
Let $\phi=(\mathrm{obj},\mathrm{lift})$ be as above.
\begin{enumerate}
  \item[\textbf{C0} (codomain).]
        For all $s$ and $h\in\Pos(D)\bigl(\mathrm{obj}(s)\bigr)$,
        \[
          \mathrm{obj}\bigl(
            \mathrm{cod}\bigl(s,\,\mathrm{lift}_s(h)\bigr)
          \bigr)
          \;=\;
          \mathrm{cod}\bigl(\mathrm{obj}(s),\,h\bigr).
        \]
  \item[\textbf{C1} (unit).]
        For all $s$,
        \[
          \mathrm{lift}_s\bigl(\id_{\mathrm{obj}(s)}\bigr)
          \;=\;
          \id_s.
        \]
  \item[\textbf{C2} (composition).]
        For all $s$, $h$, and $k$ of the appropriate types,
        \[
          \mathrm{lift}_s\bigl(k\circ h\bigr)
          \;=\;
          \mathrm{lift}_s(h)
          \;\circ\;
          \mathrm{lift}_{
            \mathrm{cod}(s,\mathrm{lift}_s(h))
          }
            \bigl((\mathrm{C0}_{s,h})^{-1}\cdot k\bigr),
        \]
        with transport of $k$ along C0 as in M2.
\end{enumerate}
\end{definition}

\begin{remark}[Variance]
A functor $\cat(C)\to\cat(D)$ would supply, for each $s$, a map
$\Pos(C)(s)\to\Pos(D)\bigl(F(s)\bigr)$---covariant on hom-sets. A cofunctor
supplies the opposite map. The object map still runs $C\to D$; only the action
on morphisms is reversed. This is the variance forced by container morphisms.
\end{remark}

%======================================================================
\section{The hom-set bijection}
%======================================================================

Under the dictionary of Definition~\ref{def:dictionary}, the data and laws of
a directed-container morphism are literally those of a cofunctor: $f$ is
$\mathrm{obj}$, $f^\sharp$ is $\mathrm{lift}$, and M0/M1/M2 are C0/C1/C2.

\begin{definition}[From morphisms to cofunctors]\label{def:to-cof}
Every directed-container morphism $\varphi=(f,f^\sharp):C\Rightarrow D$
determines a cofunctor $\varphi^\wedge:\cat(C)\to\cat(D)$ with
$\mathrm{obj}=f$ and $\mathrm{lift}=f^\sharp$.
\end{definition}

\begin{definition}[From cofunctors to morphisms]\label{def:to-dcont}
Every cofunctor $\phi=(\mathrm{obj},\mathrm{lift})$ determines a
directed-container morphism $\phi^\vee:C\Rightarrow D$ with shape map
$\mathrm{obj}$ and position map $\mathrm{lift}$.
\end{definition}

\begin{theorem}[Round-trip, morphism side]\label{thm:rt-morph}
For every directed-container morphism $\varphi$,
\begin{equation}
  \bigl(\varphi^\wedge\bigr)^\vee \;=\; \varphi.
\end{equation}
\end{theorem}

\begin{proof}
By lean.
\end{proof}

\begin{theorem}[Round-trip, cofunctor side]\label{thm:rt-cof}
For every cofunctor $\phi$,
\begin{equation}
  \bigl(\phi^\vee\bigr)^\wedge \;=\; \phi.
\end{equation}
\end{theorem}

\begin{proof}
By lean.
\end{proof}

\begin{corollary}[Hom-set bijection]\label{cor:hom-bij}
For any directed containers $C$ and $D$, there is a bijection
\begin{equation}
  \bigl\{\,\text{directed-container morphisms } C\Rightarrow D\,\bigr\}
  \;\cong\;
  \bigl\{\,\text{cofunctors } \cat(C)\to\cat(D)\,\bigr\},
\end{equation}
implemented by $(\,\cdot\,)^\wedge$ and $(\,\cdot\,)^\vee$.
\end{corollary}

\begin{proof}
By lean. (The two constructions are mutually inverse by
Theorems~\ref{thm:rt-morph} and~\ref{thm:rt-cof}.)
\end{proof}

\begin{remark}
Both round-trips are definitional: the two structures carry identical fields,
and the laws are the same equalities under renaming. No computation is
required beyond unfolding definitions.
\end{remark}

%======================================================================
\section{The category of cofunctors}
%======================================================================

\begin{definition}[Identity cofunctor]\label{def:cof-id}
The \textbf{identity cofunctor} on $\cat(C)$ has object map the identity on
shapes and lift, at each shape, the identity on positions. The laws C0--C2
hold strictly (transports along reflexivity vanish).
\end{definition}

\begin{definition}[Composition of cofunctors]\label{def:cof-comp}
Let $\phi:\cat(C)\to\cat(D)$ and $\psi:\cat(D)\to\cat(E)$ be cofunctors.
Their \textbf{composite} $\psi\circ\phi:\cat(C)\to\cat(E)$ has
\begin{enumerate}
  \item object map $\mathrm{obj}_\psi\circ\mathrm{obj}_\phi$, and
  \item lift at $s$ given by
        \[
          m \;\mapsto\;
          \mathrm{lift}_\phi\bigl(s,\,
            \mathrm{lift}_\psi\bigl(\mathrm{obj}_\phi(s),\,m\bigr)
          \bigr)
        \]
        (lifts compose backwards).
\end{enumerate}
The composite codomain law chains the two factors' C0 laws; the composite unit
law chains the two C1 laws; the composite composition law uses the factors' C2
laws together with a transport identity recorded next.
\end{definition}

\begin{lemma}[Composite transport]\label{lem:comp-transport}
Let $\psi:\cat(D)\to\cat(E)$ be a cofunctor, let $x,y\in\Shape(D)$ with
$e:y=x$, and let $z\in\Shape(E)$ with $e_2:\mathrm{obj}_\psi(x)=z$. For
$k\in\Pos(E)(z)$,
\begin{equation}
  e^{-1}\cdot\bigl(
    \mathrm{lift}_\psi\bigl(x,\, e_2^{-1}\cdot k\bigr)
  \bigr)
  \;=\;
  \mathrm{lift}_\psi\bigl(
    y,\,
    \bigl((\mathrm{obj}_\psi\cdot e)\cdot e_2\bigr)^{-1}\cdot k
  \bigr).
\end{equation}
\end{lemma}

\begin{proof}
By lean.
\end{proof}

\begin{remark}
This is the one genuinely dependent step in the composition of cofunctors:
the inner position, transported across both factors' C0 equalities in one hop,
equals the two-hop transport. After substituting the shape equalities the
identity is reflexivity.
\end{remark}

\begin{proposition}[Left unit for cofunctors]\label{prop:cof-id-comp}
For every cofunctor $\phi:\cat(C)\to\cat(D)$,
\begin{equation}
  \phi \;\circ\; \id_{\cat(C)} \;=\; \phi.
\end{equation}
\end{proposition}

\begin{proof}
By lean.
\end{proof}

\begin{proposition}[Right unit for cofunctors]\label{prop:cof-comp-id}
For every cofunctor $\phi:\cat(C)\to\cat(D)$,
\begin{equation}
  \id_{\cat(D)} \;\circ\; \phi \;=\; \phi.
\end{equation}
\end{proposition}

\begin{proof}
By lean.
\end{proof}

\begin{proposition}[Associativity for cofunctors]\label{prop:cof-assoc}
Composition of cofunctors is associative:
\begin{equation}
  \chi \;\circ\; (\psi \;\circ\; \phi)
  \;=\;
  (\chi \;\circ\; \psi) \;\circ\; \phi.
\end{equation}
\end{proposition}

\begin{proof}
By lean.
\end{proof}

\begin{theorem}[The category $\Cof$]\label{thm:cof-cat}
Directed containers and cofunctors form a strict category $\Cof$: identities
and composition as above, with unit and associativity laws holding on the
nose.
\end{theorem}

\begin{proof}
By lean. (Propositions~\ref{prop:cof-id-comp}, \ref{prop:cof-comp-id},
and~\ref{prop:cof-assoc}.)
\end{proof}

%======================================================================
\section{Isomorphism of categories}
%======================================================================

Transporting the identity cofunctor and cofunctor composition across the
hom-set bijection yields identity and composition for directed-container
morphisms.

\begin{definition}[Identity and composition in $\DCont$]\label{def:dcont-ops}
The \textbf{identity} directed-container morphism on $C$ is
$(\id_{\cat(C)})^\vee$. The \textbf{composite} of morphisms
$\varphi:C\Rightarrow D$ and $\psi:D\Rightarrow E$ is
$\bigl(\psi^\wedge\circ\varphi^\wedge\bigr)^\vee$.
\end{definition}

\begin{proposition}[Preservation of identities]\label{prop:phi-id}
The assignment $\Phi(\varphi)=\varphi^\wedge$ preserves identities:
\begin{equation}
  (\id_C)^\wedge \;=\; \id_{\cat(C)}.
\end{equation}
\end{proposition}

\begin{proof}
By lean.
\end{proof}

\begin{proposition}[Preservation of composition]\label{prop:phi-comp}
The assignment $\Phi(\varphi)=\varphi^\wedge$ preserves composition:
\begin{equation}
  (\psi\circ\varphi)^\wedge
  \;=\;
  \psi^\wedge \;\circ\; \varphi^\wedge.
\end{equation}
\end{proposition}

\begin{proof}
By lean.
\end{proof}

\begin{theorem}[Isomorphism of categories]\label{thm:iso-cats}
The assignment $\Phi$ is an isomorphism of categories
\begin{equation}
  \DCont \;\cong\; \Cof.
\end{equation}
It is the identity on objects, bijective on every hom-set
(Corollary~\ref{cor:hom-bij}), and functorial
(Propositions~\ref{prop:phi-id} and~\ref{prop:phi-comp}).
\end{theorem}

\begin{proof}
By lean.
\end{proof}

%======================================================================
\section{Discussion}
%======================================================================

\paragraph{What has been established.}
Directed-container morphisms and cofunctors of the presented categories are
the same mathematical objects under a change of vocabulary. The category they
form is strict, and the identification is an isomorphism of categories. In
particular, there are no exotic directed-container morphisms beyond cofunctors,
and no exotic cofunctors beyond directed-container morphisms.

\paragraph{Why not functors?}
The position component of a container morphism is contravariant by the
semantics of extensions: positions of the target are pulled back to positions
of the source so that labellings can be reindexed. Functors act covariantly on
hom-sets. The mismatch is definitional, not accidental; cofunctors are the
correct dual notion.

\paragraph{Scope.}
This note takes the object-level dictionary (directed containers as
categories) as given and focuses on the morphism-level identification. The
full round-trip at the object level---recovering a directed container from its
presented category up to isomorphism of positions---belongs to the companion
note on directed containers as categories.

\paragraph{Machine checking.}
Every definition and theorem statement above corresponds to a declaration in a
machine-checked formalisation. The phrase ``by lean'' indicates that the proof
has been verified by a proof assistant.

%======================================================================
\section*{Acknowledgements}
%======================================================================

The mathematical content follows the directed-container theory of Ahman,
Chapman, and Uustalu, and the cofunctor perspective associated with the
Shapiro--Spivak polynomial dictionary. The presentation here is a prose
rendering of a formal development intended for human review of the theorem
statements.

\begin{thebibliography}{9}

\bibitem{ahman2012}
D.~Ahman, J.~Chapman, and T.~Uustalu.
\newblock When is a container a comonad?
\newblock In \emph{FoSSaCS}, volume~7213 of \emph{LNCS}, pages 74--88.
\newblock Springer, 2012.

\bibitem{ahman2016}
D.~Ahman and T.~Uustalu.
\newblock Directed containers as categories.
\newblock In \emph{MSFP}, volume~207 of \emph{EPTCS}, pages 89--98, 2016.

\bibitem{shapiro2023}
B.~Shapiro and D.~I.~Spivak.
\newblock Structures on categories of polynomials.
\newblock Preprint, 2023.

\end{thebibliography}

\end{document}
